\chapter{Prakata}

Buku ini ditulis untuk pembaca yang sudah dapat memprogram dan kini perlu menjalankan, mengamati, mengotomasi, dan akhirnya membangkitkan sistem perangkat lunak dari model. Pembahasannya dimulai dari sistem yang benar-benar berjalan, lalu bergerak ke pipeline otomasi, dan berakhir pada model yang menghasilkan artefak operasional.

Pembaca sasarannya adalah mahasiswa tingkat akhir dan pascasarjana bidang informatika, serta praktisi yang bekerja di sekitar DevOps dan rekayasa berbasis model. Materi ini dipakai pada mata kuliah Advanced Software Engineering and DevOps di Program Studi Informatika, Universitas Pradita, dengan empat belas pertemuan.

Seluruh studi kasus dibangun dan dijalankan oleh penulis. Pipeline pemrosesan video, \textit{stack} layanan \textit{text to speech}, metamodel MediaFlow, DSL tekstual dan visual, transformasi, generator, aturan validasi, serta sistem multi-agen tersedia terbuka pada \url{https://github.com/alfa-yohannis/advanced-software-engineering}. Kode tiap bab berada di folder \texttt{projects/session-NN} sesuai nomor babnya, dan naskah buku ini ada di folder \texttt{module}. Perintah dan keluaran yang tercetak di buku ini berasal dari hasil menjalankannya.

Tiap bab dibuka dengan tujuan pembelajaran dan ditutup dengan ringkasan serta latihan. Latihan dirancang untuk dikerjakan di komputer pembaca, dan sebagian menuntut pembaca merusak sistem lebih dahulu agar dapat mengamati cara sistem itu gagal.

Penulis berterima kasih kepada mahasiswa yang mengikuti mata kuliah ini dan memberi banyak koreksi melalui pertanyaan mereka di kelas. Saran dan koreksi dari pembaca dapat disampaikan melalui penerbit.

\bigskip
\begin{flushright}
Tangerang, \isian{bulan dan tahun terbit}\\[4pt]
Penulis
\end{flushright}
